% संपूर्ण मराठी ग्रंथातील प्रकरण 20: अंकगणिताची प्रतिमाने
% हा संपादनयोग्य स्रोतखंड आहे. संकलनासाठी संपूर्ण स्रोतसंग्रहातील
% अक्षररूपे, पूर्वभाग, चिन्हव्याख्या आणि आकृती-स्रोत आवश्यक आहेत.
% मूळ: Open Logic Project; मजकूर आणि रूपांतर CC BY 4.0.
% यंत्रानुवाद, दुरुस्ती आणि तपासणी: OpenAI Codex —
% GPT-5.6 Sol आणि GPT-6 Sol, दोन्ही Ultra effort.
% दुरुस्ती-पुनरावलोकन: GPT-6.1 Sol, Ultra effort.
\chapter{अंकगणिताची प्रतिमाने}\label{mod:mar::chap}








\section{परिचय}\label{mod:mar:int:sec}

अंकगणिताचे \emph{मानक प्रतिमान} म्हणजे अशी संरचना~$\Struct{N}$ की
$\Domain{N} = \Nat$ आणि तिच्यात $\Obj{0}$, $\prime$, $+$, $\times$ व
$<$ यांचे अपेक्षित नेहमीच्या अर्थाने अर्थनिर्धारण केलेले असते. म्हणजे,
$\Obj{0}$ हे $0$ असते, $\prime$ हे उत्तरवर्ती फलन असते, $+$ चा अर्थ
बेरीज आणि $\times$ चा अर्थ~$\Nat$ मधील संख्यांचा गुणाकार असा असतो.
अधिक स्पष्टपणे,
\begin{align*}
  \Assign{\Obj{0}}{N} & = 0\\
  \Assign{\prime}{N}(n) & = n + 1\\
  \Assign{+}{N}(n, m) & = n + m\\
  \Assign{\times}{N}(n, m) & = nm
\end{align*}
अर्थात, $\Lang{L_A}$ साठी अशाही रचना असतात की त्यांचे प्रांत~$\Nat$
नसतात. उदाहरणार्थ, $\Domain{M} = \{a\}^*$ (एकाच $a$ या चिन्हाच्या
सांत क्रमिका, म्हणजे $\emptyset$, $a$, $aa$, $aaa$, \dots) असा प्रांत
असलेली $\Struct{M}$ घेता येते आणि पुढील अर्थनिर्धारण करता येते:
\begin{align*}
  \Assign{\Obj{0}}{M} & = \emptyset\\
  \Assign{\prime}{M}(s) & = s \concat a\\
  \Assign{+}{M}(n, m) & = a^{n + m}\\
  \Assign{\times}{M}(n, m) & = a^{nm}
\end{align*}
या दोन रचना ``मूलतः एकच'' आहेत: फरक फक्त प्रांताचे
घटक कोणते आहेत यात आहे; अर्थनिर्धारण फलनांमुळे प्रांताचे
घटक परस्पर कसे संबंधित आहेत यात नाही. अशा दोन संरचना ना
\emph{समरूप} म्हणतात.

संहतता प्रमेयावरून सहज निष्पन्न होते की~$\Struct{N}$ मध्ये सत्य असलेल्या
कोणत्याही उपपत्तीला~$\Struct{N}$ शी समरूप नसलेली प्रतिमानेही असतात. अशा
रचनांना \emph{अमानक} म्हणतात. त्यांच्याबद्दलची रंजक गोष्ट अशी की मानक
प्रतिमानातील घटक (म्हणजे $\Struct{N}$ मधील, तसेच तिच्याशी समरूप
असलेल्या सर्व संरचना मधील घटक) मानक संख्यांक~$\num{n}$ यांच्या
मूल्यांनी पूर्णपणे व्यापलेले असतात, म्हणजे
\[\Domain{N} = \Setabs{\Value{\num{n}}{N}}{n \in \Nat}
\]
पण अमानक प्रतिमानांत तसे नसते: $\Struct{M}$ अमानक असेल, तर किमान एक
$x \in \Domain{M}$ असा असतो की प्रत्येक~$n$ साठी
$x \neq \Value{\num{n}}{M}$.

हे अमानक घटक फार रंजक आहेत: ते ``अनंत नैसर्गिक संख्या'' आहेत. पण त्यांच्या
अस्तित्वामुळे अपूर्णता-घटनांचे एका अर्थी स्पष्टीकरणही मिळते. उदाहरणार्थ,
पेआनो अंकगणितासाठीचे सुसंगतता-विधान $\OCon[\Th{PA}]$, म्हणजे $\lnot
\lexists[x][\OPrf[\Th{PA}](x, \gn{\lfalse})]$, विचारात घ्या.
$\Th{PA}$ कडून $\OCon[\Th{PA}]$ किंवा $\lnot \OCon[\Th{PA}]$ यांपैकी
काहीही सिद्ध होत नसल्यामुळे यांपैकी कोणतेही विधान~$\Th{PA}$ मध्ये सुसंगतपणे
जोडता येते. $\Th{PA}$ सुसंगत असल्यामुळे
$\Sat{N}{\OCon[\Th{PA}]}$, आणि म्हणून $\Sat/{N}{\lnot
  \OCon[\Th{PA}]}$. त्यामुळे $\Struct{N}$ हे $\Th{PA}
\cup \{\lnot \OCon[\Th{PA}]\}$ चे \emph{प्रतिमान नाही}, आणि त्या उपपत्तीची
सर्व प्रतिमाने अमानकच असली पाहिजेत. $\Th{PA} \cup \{\lnot
\OCon[\Th{PA}]\}$ च्या प्रतिमानात असा एखादा घटक असलाच पाहिजे की
तो $\lexists[x][\OPrf[\Th{PA}](x, \gn{\lfalse})]$ सत्य करणारा साक्षी
ठरेल, म्हणजे~$\Th{PA}$ पासून व्याघाताच्या निष्पत्तीची गोडेल संख्या.

असा घटक मानक असू शकत नाही---कारण प्रत्येक~$n$ साठी
$\Th{PA} \Proves \lnot \OPrf[\Th{PA}](\num{n}, \gn{\lfalse})$.











\section{अंकगणिताची मानक प्रतिमाने}\label{mod:mar:stm:sec}

अंकगणिताची भाषा~$\Lang{L_A}$ ही अर्थातच संख्यांविषयी, विशेषतः नैसर्गिक
संख्यांविषयी, अभिप्रेत आहे. म्हणून ``ते'' मानक प्रतिमान~$\Struct{N}$ विशेष
आहे: आपल्याला ज्याविषयी बोलायचे आहे तेच हे प्रतिमान. पण तर्कशास्त्रात
आपल्याला अनेकदा फक्त रचनात्मक गुणधर्मांत रस असतो, आणि समरूप असलेल्या कोणत्याही
दोन संरचनांमध्ये ते समान असतात. त्यामुळे व्याप्ती थोडी वाढवून
$\Struct{N}$ शी समरूप असलेली कोणतीही संरचना ``मानक'' मानता येते.

\begin{defn}
$\Lang{L_A}$ साठीची एखादी संरचना जर~$\Struct{N}$ शी समरूप असेल,
तर ती \emph{मानक} आहे.
\end{defn}

\begin{prop}
\label{mod:mar:stm:prop:standard-domain} एखादी संरचना~$\Struct{M}$ मानक असेल,
तर तिचा प्रांत म्हणजे मानक संख्यांकांच्या मूल्यांचा संच असतो, म्हणजे
\[\Domain{M} = \Setabs{\Value{\num{n}}{M}}{n \in \Nat}
\]
\end{prop}

\begin{proof}
प्रत्येक $\Value{\num{n}}{M} \in \Domain{M}$ हे स्पष्ट आहे. प्रत्येक
$x \in \Domain{M}$ हा कोणत्यातरी~$n$ साठी $\Value{\num{n}}{M}$ बरोबर आहे,
एवढेच दाखवायचे आहे. $\Struct{M}$ मानक असल्यामुळे ते~$\Struct{N}$ शी समरूप
आहे. $g\colon \Nat \to \Domain{M}$ हे समरूपण आहे असे समजा. मग
$g(n) = g(\Value{\num{n}}{N}) = \Value{\num{n}}{M}$. पण प्रत्येक
$x \in \Domain{M}$ साठी $g(n) = x$ असे एक~$n \in \Nat$ असते, कारण
$g$ हे आच्छादक आहे.
\end{proof}

\begin{explain}
$\Lang{L_A}$ साठीची एखादी रचना~$\Struct{M}$ मानक असेल, तर तिच्या
प्रांत मधील सर्व घटकांना $\num{0}$, $\num{1}$, $\num{2}$, \dots हे
मानक संख्यांक, म्हणजेच $\Obj{0}$, $\Obj{0}'$, $\Obj{0}''$ इत्यादी पदे,
नावे देऊ शकतात. अर्थात, याचा अर्थ $\Domain{M}$ मधील घटक
\emph{संख्याच आहेत} असा नाही; $\Domain{N}$ मधील संख्या जशा निवडून दाखवता
येतात, त्याच रीतीने हे घटक निवडून दाखवता येतात एवढाच त्याचा अर्थ आहे.
\end{explain}

\begin{prob}
\ref{mod:mar:stm:prop:standard-domain} चा व्यत्यास असत्य आहे हे दाखवा;
म्हणजे, $\Domain{M} = \Setabs{\Value{\num{n}}{M}}{n \in \Nat}$ असूनही
$\Struct{N}$ शी समरूप नसलेल्या एखाद्या संरचना~$\Struct{M}$ चे
उदाहरण द्या.
\end{prob}

\begin{prop}
\label{mod:mar:stm:prop:thq-standard}
जर $\Sat{M}{\Th{Q}}$ आणि $\Domain{M} =
\Setabs{\Value{\num{n}}{M}}{n \in \Nat}$, तर $\Struct{M}$ मानक आहे.
\end{prop}

\begin{proof}
$\Struct{M}$ हे~$\Struct{N}$ शी समरूप आहे हे दाखवायचे आहे.
$g(n) = \Value{\num{n}}{M}$ अशी व्याख्या असलेले
$g\colon \Nat \to \Domain{M}$ फलन विचारात घ्या. गृहीतकानुसार $g$ हे
आच्छादक आहे. ते एकास-एक सुद्धा आहे: $n \neq m$ असेल तेव्हा
$\Th{Q} \Proves \eq/[\num{n}][\num{m}]$. म्हणून
$\Sat{M}{\Th{Q}}$ असल्यामुळे, $n \neq m$ असेल तेव्हा
$\Sat{M}{\eq/[\num{n}][\num{m}]}$. त्यामुळे $n \neq m$ असल्यास
$\Value{\num{n}}{M} \neq \Value{\num{m}}{M}$, म्हणजे $g(n) \neq g(m)$.

$g$ हे समरूपण आहे हेही पडताळायचे आहे.
\begin{enumerate}
\item $\Assign{\Obj{0}}{N} = 0$ असल्यामुळे
  $g(\Assign{\Obj{0}}{N}) = g(0)$. $g$ च्या व्याख्येवरून
  $g(0) = \Value{\num{0}}{M}$. पण $\num{0}$ म्हणजे फक्त $\Obj{0}$;
  आणि एखादे पद स्थिरांक असेल, तर त्याचे मूल्य त्या संरचना ने
  त्या स्थिरांक ला नेमून दिलेले मूल्य असते, म्हणजे
  $\Value{\Obj{0}}{M} = \Assign{\Obj{0}}{M}$. म्हणून अपेक्षेप्रमाणे
  $g(\Assign{\Obj{0}}{N}) = \Assign{\Obj{0}}{M}$.
\item $\Struct{N}$ मध्ये $\prime$ हे~$\Nat$ वरील उत्तरवर्ती फलन असल्यामुळे
  $g(\Assign{\prime}{N}(n)) = g(n+1)$. पुढे, $g$ च्या व्याख्येवरून
  $g(n+1) = \Value{\num{n+1}}{M}$. पण $\num{n+1}$ आणि $\num{n}'$ हे
  एकच पद आहे, म्हणून $\Value{\num{n+1}}{M} = \Value{\num{n}'}{M}$.
  मूल्य-फलनाच्या व्याख्येवरून हे
  $= \Assign{\prime}{M}(\Value{\num{n}}{M})$. तसेच
  $\Value{\num{n}}{M} = g(n)$ असल्यामुळे
  $g(\Assign{\prime}{N}(n)) = \Assign{\prime}{M}(g(n))$ मिळते.
\item $\Struct{N}$ मध्ये $+$ हे~$\Nat$ वरील बेरीज-फलन असल्यामुळे
  $g(\Assign{+}{N}(n,m)) = g(n+m)$. पुढे, $g$ च्या व्याख्येवरून
  $g(n+m) = \Value{\num{n+m}}{M}$. पण $\Th{Q} \Proves
  \num{n+m} = (\num{n} + \num{m})$, म्हणून $\Value{\num{n+m}}{M} =
  \Value{\num{n}+\num{m}}{M}$. मूल्य-फलनाच्या व्याख्येवरून हे $=
  \Assign{+}{M}(\Value{\num{n}}{M},\Value{\num{m}}{M})$. तसेच
  $\Value{\num{n}}{M} = g(n)$ आणि $\Value{\num{m}}{M} = g(m)$ असल्यामुळे
  $g(\Assign{+}{N}(n, m)) = \Assign{+}{M}(g(n), g(m))$ मिळते.
\item $g(\Assign{\times}{N}(n, m)) = \Assign{\times}{M}(g(n), g(m))$:
  सराव.
\item $\tuple{n,m} \in \Assign{<}{N}$ हे तेव्हाच, जेव्हा $n < m$. जर
  $n < m$, तर $\Th{Q} \Proves \num{n} < \num{m}$ आणि
  $\Sat{M}{\num{n} < \num{m}}$ सुद्धा. त्यामुळे
  $\tuple{\Value{\num{n}}{M}, \Value{\num{m}}{M}} \in \Assign{<}{M}$,
  म्हणजे $\tuple{g(n), g(m)} \in \Assign{<}{M}$. जर $n \not< m$, तर
  $\Th{Q} \Proves \lnot \num{n} < \num{m}$ आणि परिणामी
  $\Sat/{M}{\num{n} < \num{m}}$. त्यामुळे पूर्वीप्रमाणे
  $\tuple{g(n), g(m)} \notin \Assign{<}{M}$. दोन्ही एकत्र केल्यास:
  $\tuple{n,m} \in \Assign{<}{N}$ हे तेव्हाच, जेव्हा
  $\tuple{g(n), g(m)} \in \Assign{<}{M}$.
\end{enumerate}
\end{proof}

\begin{explain}
$\Lang{L_A}$ साठीच्या इतर कोणत्याही संरचना~$\Struct{M}$ च्या
प्रांतात~$\Nat$ कडून संगती ठरवण्याचा सर्वांत साहजिक मार्ग म्हणजे फलन~$g$;
कारण अशा प्रत्येक $\Struct{M}$ मध्ये $\num{0}$, $\num{1}$, $\num{2}$
इत्यादींनी नाव दिलेले घटक असतात. त्यामुळे, $\Struct{M}$ ने
$\num{n}$ विषयीची किमान काही मूलभूत विधाने~$\Struct{N}$ प्रमाणेच सत्य
केली आणि $g$ हे एकास-एक व आच्छादकही असेल, तर $g$ हे समरूपण ठरेल यात नवल
नाही. किंबहुना, $\Domain{M}$ मध्ये $\num{n}$ यांनी नाव दिलेल्यांव्यतिरिक्त
दुसरे घटक नसतील, तर असे समरूपण एकमेव असते.
\end{explain}

\begin{prop}
\label{mod:mar:stm:prop:thq-unique-iso} जर $\Struct{M}$ मानक असेल, तर
\ref{mod:mar:stm:prop:thq-standard} च्या पुराव्यातील $g$ हे~$\Struct{N}$ कडून
$\Struct{M}$ कडे जाणारे एकमेव समरूपण आहे.
\end{prop}

\begin{proof}
$h\colon \Nat \to \Domain{M}$ हे $\Struct{N}$ आणि~$\Struct{M}$ यांमधील
समरूपण आहे असे समजा. $n$ वरील विगमनाने $g = h$ हे दाखवू. $n = 0$ असेल,
तर $g$ च्या व्याख्येवरून $g(0) = \Assign{\Obj{0}}{M}$. पण $h$ हे समरूपण
असल्यामुळे $h(0) = h(\Assign{\Obj{0}}{N}) =\Assign{\Obj{0}}{M}$;
म्हणून $g(0) = h(0)$.

आता $n+1$ चे प्रकरण विचारात घ्या. आपल्याला
\begin{align*}
  g(n+1) & = \Value{\num{n+1}}{M} \text{ हे~$g$ च्या व्याख्येवरून}\\
  & = \Value{\num{n}'}{M} \text{ कारण $\num{n+1}\ident \num{n}'$}\\
  & = \Assign{\prime}{M}(\Value{\num{n}}{M})
    \text{ हे $\Value{t'}{M}$ च्या व्याख्येवरून}\\
  & = \Assign{\prime}{M}(g(n))  \text{ हे~$g$ च्या व्याख्येवरून}\\
  & = \Assign{\prime}{M}(h(n)) \text{ हे विगमन गृहीतकावरून}\\
  & = h(\Assign{\prime}{N}(n)) \text{ कारण $h$ हे समरूपण आहे}\\
  & = h(n+1)
\end{align*}
मिळते.
\end{proof}

\begin{explain}
कोणत्याही गणनीय अनंत संच~$M$ साठी $\Nat$ आणि $M$ यांमध्ये
एकास-एक आच्छादन असते; म्हणून असा प्रत्येक संच~$M$ हा एखाद्या मानक
प्रतिमान~$\Struct{M}$ चा संभाव्य प्रांत आहे. किंबहुना, एकदा $z \in M$
ही वस्तू आणि योग्य फलन $s$ ही अनुक्रमे $\Assign{\Obj{0}}{M}$ आणि
$\Assign{\prime}{M}$ म्हणून निवडली, की $+$, $\times$ आणि $<$ यांचे
अर्थनिर्धारण आधीच निश्चित होते. $s\colon M \to M \setminus \{z\}$ अशी
एकास-एक आणि आच्छादक दोन्ही असलेली फलनेच मानक प्रतिमानात
$\Assign{\prime}{M}$ म्हणून योग्य असतात. $s$ च्या व्याप्तीत~$z$ असू शकत
नाही; अन्यथा $\lforall[x][\eq/[\Obj 0][x']]$ असत्य होईल. हे वाक्य
$\Struct{N}$ मध्ये सत्य आहे, म्हणून $\Struct{M}$ नेही ते सत्य केले पाहिजे.
फलन~$s$ हे एकास-एक असले पाहिजे, कारण $\Struct{N}$ मधील उत्तरवर्ती
फलन~$\Assign{\prime}{N}$ तसे आहे; आणि $\Assign{\prime}{N}$ हे
एकास-एक आहे हे~$\Struct{N}$ मध्ये सत्य असलेल्या वाक्य ने
व्यक्त होते. ते आच्छादक सुद्धा असले पाहिजे; अन्यथा असा एखादा
$x \in M \setminus \{z\}$ असेल की तो~$s$ च्या व्याप्तीत नसेल, म्हणजे
वाक्य $\lforall[x][(\eq[x][\Obj 0] \lor
\lexists[y][\eq[y'][x]])]$ हे~$\Struct{M}$ मध्ये असत्य होईल---पण ते
$\Struct{N}$ मध्ये सत्य आहे.

\end{explain}












\section{अमानक प्रतिमाने}\label{mod:mar:nst:sec}

\begin{explain}
$\Lang{L_A}$ साठीची एखादी संरचना~$\Struct{N}$ शी समरूप असेल,
तर तिला मानक म्हणतो. एखादी संरचना~$\Struct{N}$ शी समरूप नसेल,
तर तिला अमानक म्हणतो.
\end{explain}

\begin{defn}
$\Lang{L_A}$ साठीची संरचना~$\Struct{M}$ जर~$\Struct{N}$ शी समरूप
नसेल, तर ती \emph{अमानक} आहे. प्रांतातील जे घटक
$x \in \Domain{M}$ कोणत्यातरी $n \in \Nat$ साठी
$\Value{\num{n}}{M}$ बरोबर आहेत, त्यांना ($\Struct{M}$ मधील)
\emph{मानक संख्या} म्हणतात; आणि जे तसे नाहीत त्यांना \emph{अमानक संख्या}
म्हणतात.
\end{defn}

\begin{explain}
\ref{mod:mar:stm:prop:standard-domain} नुसार, $\Lang{L_A}$ साठीच्या कोणत्याही
मानक संरचनांमध्ये फक्त मानक घटक असतात. त्यामुळे अमानक
संरचनांमध्ये किमान एक अमानक घटक असलाच पाहिजे. किंबहुना, अमानक
घटक अस्तित्वात असणे हे ती संरचना अमानक असल्याची हमी देते.
\end{explain}

\begin{prop}
$\Lang{L_A}$ साठीच्या संरचना~$\Struct{M}$ मध्ये अमानक संख्या
असेल, तर $\Struct{M}$ अमानक आहे.
\end{prop}

\begin{proof}
तसे नाही असे समजा; म्हणजे $\Struct{M}$ मानक असूनही तिच्यात अमानक संख्या
$x$ आहे असे समजा. $g\colon \Nat \to \Domain{M}$ हे समरूपण असू द्या.
$n$ वरील विगमनाने $g(\Value{\num{n}}{N}) = \Value{\num{n}}{M}$ हे सहज
दिसते. दुसऱ्या शब्दांत, $g$ हे~$\Struct{N}$ मधील मानक संख्यांना
$\Struct{M}$ मधील मानक संख्यांवर पाठवते. $\Struct{M}$ मध्ये अमानक संख्या
असेल, तर गृहीतकाच्या विरुद्ध $g$ हे आच्छादक असू शकत नाही.
\end{proof}

\begin{prob}
$\Th{Q}$ मध्ये पुढील स्वयंसिद्धके आहेत हे आठवा:
\begin{align*}
& \lforall[x][\lforall[y][(\eq[x'][y'] \lif \eq[x][y])]] \tag{$\mathsf{Q}_1$}\\
& \lforall[x][\eq/[\Obj 0][x']] \tag{$\mathsf{Q}_2$}\\
& \lforall[x][(\eq[x][\Obj 0] \lor \lexists[y][\eq[x][y']])] \tag{$\mathsf{Q}_3$}
\end{align*}
पुढील अटी पूर्ण करणाऱ्या संरचना~$\Struct{M_1}$, $\Struct{M_2}$,
$\Struct{M_3}$ द्या:
\begin{enumerate}
\item $\Sat{M_1}{\mathsf{Q}_1}$, $\Sat{M_1}{\mathsf{Q}_2}$, $\Sat/{M_1}{\mathsf{Q}_3}$;
\item $\Sat{M_2}{\mathsf{Q}_1}$, $\Sat/{M_2}{\mathsf{Q}_2}$, $\Sat{M_2}{\mathsf{Q}_3}$; आणि
\item $\Sat/{M_3}{\mathsf{Q}_1}$, $\Sat{M_3}{\mathsf{Q}_2}$, $\Sat{M_3}{\mathsf{Q}_3}$.
\end{enumerate}
अर्थात, प्रत्येकासाठी फक्त~$\Assign{\Obj{0}}{M_i}$ आणि
$\Assign{\prime}{M_i}$ निश्चित करायची आहेत.
\end{prob}

\begin{explain}
$\Lang{L_A}$ साठी अमानक संरचना निश्चित करणे पुरेसे सोपे आहे.
उदाहरणार्थ, प्रांत~$\Int$ असलेली रचना घ्या आणि सर्व अतार्किक चिन्हांचे
नेहमीप्रमाणे अर्थनिर्धारण करा. ऋण संख्या कोणत्याही~$n$ साठी $\num{n}$ ची
मूल्ये नसल्यामुळे ही रचना अमानक आहे. अर्थात, $\Struct{N}$ प्रमाणेच ती
वाक्ये सत्य करते या अर्थाने ती अंकगणिताचे \emph{प्रतिमान} नसेल. उदाहरणार्थ,
$\lforall[x][\eq/[x'][\Obj{0}]]$ हे असत्य आहे. तरीसुद्धा, संहतता प्रमेय
वापरून अंकगणिताची अमानक प्रतिमाने अस्तित्वात आहेत हे सहज सिद्ध करता येते.
\end{explain}

\begin{prop}
$\Th{TA} = \Setabs{\varphi}{\Sat{N}{\varphi}}$ ही~$\Struct{N}$ ची उपपत्ती असू द्या.
$\Th{TA}$ ला गणनीय अमानक प्रतिमान आहे.
\end{prop}

\begin{proof}
$\Lang{L_A}$ मध्ये नवीन स्थिरांक~$c$ घालून तिचा विस्तार करा आणि
पुढील वाक्यांचा संच विचारात घ्या:
\[\Gamma = \Th{TA} \cup \{\eq/[c][\num{0}], \eq/[c][\num{1}],
\eq/[c][\num{2}], \dots\}
\]
$\Gamma$ च्या कोणत्याही प्रतिमान~$\Struct{M^c}$ मध्ये
$x = \Assign{c}{M^c}$ असा अमानक घटक असेल, कारण प्रत्येक
$n \in \Nat$ साठी $x \neq \Value{\num{n}}{M}$. तसेच
$\Th{TA} \subseteq \Gamma$ असल्यामुळे अर्थातच $\Sat{M^c}{\Th{TA}}$.
फक्त~$c$ विसरून $\Struct{M^c}$ पासून $\Lang{L_A}$ साठीची
संरचना~$\Struct{M}$ बनवली, तरी तिच्या प्रांतात तो अमानक~$x$ राहतो
आणि $\Sat{M}{\Th{TA}}$ सुद्धा. $\Th{TA}$ मध्ये $c$ येत नसल्यामुळे या
दुसऱ्या दाव्याची हमी मिळते. म्हणून $\Gamma$ ला प्रतिमान आहे एवढे दाखवणे
पुरेसे आहे.


$\Gamma$ ला प्रतिमान आहे हे दाखवण्यासाठी संहतता प्रमेय वापरू. $\Gamma$ चा
प्रत्येक सांत उपसंच पूर्ततायोग्य असेल, तर $\Gamma$ सुद्धा पूर्ततायोग्य आहे.
कोणताही सांत उपसंच $\Gamma_0 \subseteq \Gamma$ विचारात घ्या.
$\Gamma_0$ मध्ये~$\Th{TA}$ मधील काही वाक्ये आणि
$\eq/[c][\num{n}]$ या रूपातील सांत वाक्ये येतात. या दुसऱ्या रूपातील एकही
वाक्य नसेल, तर c ला शून्य मूल्य देणारा N चा विस्तार त्या उपसंचाला थेट
पूर्ण करतो; म्हणून पुढे किमान एक असे वाक्य आहे असे समजा.

$\eq/[c][\num{k}] \in \Gamma_0$ असा $k$ हा सर्वांत मोठा अंक समजा.
$\Assign{c}{N_k} = k+1$ हे अर्थनिर्धारण समाविष्ट करण्यासाठी~$\Struct{N}$
चा विस्तार करून $\Struct{N_k}$ ची व्याख्या करा. $\Sat{N_k}{\Gamma_0}$:
जर $\varphi \in \Th{TA}$, तर $\Struct{N_k}$ ही~$c$ वगळता प्रत्येक बाबतीत
$\Struct{N}$ सारखीच असल्यामुळे आणि~$\varphi$ मध्ये $c$ येत नसल्यामुळे
$\Sat{N_k}{\varphi}$. तसेच $n \le k$ आणि $\Value{c}{N_k} = k+1$ असल्यामुळे
$\Sat{N_k}{\eq/[c][\num{n}]}$. म्हणून $\Gamma$ चा प्रत्येक सांत उपसंच
पूर्ततायोग्य आहे आणि संहततेनुसार संपूर्ण संचाला प्रतिमान आहे. अधोगामी
लोव्हेनहाइम--स्कोलेम प्रमेयानुसार हे प्रतिमान गणनीय घेता येते.

\end{proof}











\section{$\Th{Q}$ ची प्रतिमाने}\label{mod:mar:mdq:sec}

\begin{explain}
$\Struct{N}$ जी वाक्ये सत्य करते तीच सत्य करणाऱ्या अमानक
संरचना आहेत, म्हणजे त्या~$\Th{TA}$ ची प्रतिमाने आहेत, हे आपल्याला
माहीत आहे. $\Sat{N}{\Th{Q}}$ असल्यामुळे $\Th{TA}$ चे कोणतेही प्रतिमान
$\Th{Q}$ चेही प्रतिमान असते. $\Th{Q}$ ही~$\Th{TA}$ पेक्षा फारच दुर्बल
आहे; उदाहरणार्थ, $\Th{Q} \Proves/ \lforall[x][\lforall[y][\eq[(x +
      y)][(y+x)]]]$. दुर्बलतर उपपत्ती पूर्ण करणे सोपे असते: त्यांना अधिक
प्रतिमाने असतात. उदाहरणार्थ, $\Th{Q}$ च्या काही प्रतिमानांत
$\lforall[x][\lforall[y][\eq[(x + y)][(y+x)]]]$ असत्य असते; पण ती
$\Th{TA}$ ची किंवा त्या दृष्टीने $\Th{PA}$ चीसुद्धा प्रतिमाने असू शकत
नाहीत. $\Th{Q}$ ची प्रतिमाने तुलनेने सोपीही असतात: ती स्पष्टपणे निश्चित
करता येतात.
\end{explain}

\begin{ex}
\label{mod:mar:mdq:ex:model-K-of-Q}
$\Domain{K} = \Nat \cup \{a\}$ असा प्रांत आणि पुढील अर्थनिर्धारणे असलेली
संरचना~$\Struct{K}$ विचारात घ्या:
\begin{align*}
  \Assign{\Obj{0}}{K} & = 0\\
  \Assign{\prime}{K}(x) & =
  \begin{cases}
    x+1 & \text{जर $x\in \Nat$}\\
    a & \text{जर $x = a$}
  \end{cases}\\
  \Assign{+}{K}(x, y) & =
  \begin{cases}
    x+y & \text{जर $x$, $y \in\Nat$}\\
    a & \text{इतर प्रसंगी}
  \end{cases}\\
  \Assign{\times}{K}(x, y) & =
  \begin{cases}
    xy & \text{जर $x$, $y \in\Nat$}\\
    0 & \text{जर $x = 0$ किंवा $y = 0$}\\
    a & \text{इतर प्रसंगी}\\
  \end{cases}\\
  \Assign{<}{K} & =
  \Setabs{\tuple{x,y}}{x, y \in \Nat \text{ आणि } x<y} \cup
  \Setabs{\tuple{x,a}}{x \in \Domain{K}}
\end{align*}
$\Sat{K}{\Th{Q}}$ हे दाखवण्यासाठी $\Th{Q}$ ची सर्व स्वयंसिद्धके
$\Struct{K}$ मध्ये सत्य आहेत हे पडताळावे लागेल. सोयीसाठी,
$\Assign{\prime}{K}(x)$ करिता $x^\nssucc$ ($\Struct{K}$ मधील $x$ चा
``उत्तरवर्ती''), $\Assign{+}{K}(x, y)$ करिता $x \nsplus y$
($\Struct{K}$ मधील $x$ आणि $y$ यांची ``बेरीज''),
$\Assign{\times}{K}(x, y)$ करिता $x \nstimes y$ ($\Struct{K}$ मधील
$x$ आणि~$y$ यांचा ``गुणाकार''), आणि $\tuple{x,y} \in \Assign{<}{K}$
करिता $x \nsless y$ लिहू. या संक्षेपांनी $\Struct{K}$ मधील क्रिया अधिक
स्पष्टपणे अशा देता येतात:
\[\begin{array}{c|c}
  x & x^\nssucc \\
  \hline
  n & n+1 \\
  a & a
\end{array}
\qquad
\begin{array}{c|ccc}
  x \nsplus y & 0 & m & a \\
  \hline
  0 & 0 & m & a \\
  n & n & n+m & a \\
  a & a & a & a \\
\end{array}
\qquad
\begin{array}{c|ccc}
  x \nstimes y & 0 & m & a \\
  \hline
  0 & 0 & 0 & 0 \\
  n & 0 & nm & a \\
  a & 0 & a & a \\
\end{array}
\]
$n$, $m \in \Nat$ यांच्यासाठी $n \nsless m$ हे तेव्हाच, जेव्हा $n<m$;
आणि प्रत्येक~$x \in \Domain{K}$ साठी $x \nsless a$.

$\nssucc$ हे एकास-एक असल्यामुळे
$\Sat{K}{\lforall[x][\lforall[y][(\eq[x'][y'] \lif \eq[x][y])]]}$.
$\Struct{K}$ मध्ये $0$ हा $\nssucc$-उत्तरवर्ती नसल्यामुळे
$\Sat{K}{\lforall[x][\eq/[\Obj 0][x']]}$. प्रत्येक $n>0$ साठी
$n = (n-1)^\nssucc$ आणि $a = a^\nssucc$ असल्यामुळे
$\Sat{K}{\lforall[x][(\eq[x][\Obj 0] \lor \lexists[y][\eq[x][y']])]}$.

$n \nsplus 0 = n+0 = n$ आणि~$\nsplus$ च्या व्याख्येवरून
$a\nsplus 0 = a$ असल्यामुळे $\Sat{K}{\lforall[x][\eq[(x + \Obj 0)][x]]}$.
$\Sat{K}{\lforall[x][\lforall[y][\eq[(x + y')][(x + y)']]]}$ हे थोडे
अधिक गुंतागुंतीचे आहे. $n$, $m$ या दोन्ही मानक असतील, तर:
\begin{align*}
(n \nsplus m^\nssucc) & = (n+(m+1)) = (n+m)+1 = (n \nsplus m)^\nssucc
\intertext{कारण मानक संख्यांवर $\nsplus$ आणि $^\nssucc$ ही अनुक्रमे $+$
  आणि $\prime$ यांच्याशी जुळतात. आता $x \in \Domain{K}$ असे समजा. मग}
(x \nsplus a^\nssucc) & = (x \nsplus a) = a = a^\nssucc = (x \nsplus a)^\nssucc
\intertext{उरलेले प्रकरण $y \in \Domain{K}$ पण $x =
  a$ असण्याचे आहे. येथे $y = n$ मानक आहे की $y = a$, यानुसारही प्रकरणे
  वेगळी करावी लागतात:}
(a \nsplus n^\nssucc) & = (a \nsplus (n+1)) = a = a^\nssucc = (a \nsplus n)^\nssucc\\
(a \nsplus a^\nssucc) & = (a \nsplus a) = a = a^\nssucc = (a \nsplus a)^\nssucc
\end{align*}

अर्थात, हा तपशील आवश्यकतेपेक्षा थोडा अधिक आहे. उदाहरणार्थ,
$z$ काहीही असले तरी $a \nsplus z = a$ असल्यामुळे
$a \nsplus a^\nssucc = a$ हे लगेच निष्पन्न होते. उरलेली स्वयंसिद्धकेही
याच रीतीने पडताळता येतात.

त्यामुळे $\Struct{K}$ हे~$\Th{Q}$ चे प्रतिमान आहे. तिची ``बेरीज''~$\nsplus$
ही क्रमनिरपेक्षसुद्धा आहे. पण $\Struct{N}$ मध्ये सत्य आणि $\Struct{K}$
मध्ये असत्य अशी इतर वाक्ये आहेत, तसेच याच्या उलटही आहे. उदाहरणार्थ,
$a \nsless a$, म्हणून $\Sat{K}{\lexists[x][x < x]}$ आणि
$\Sat/{K}{\lforall[x][\lnot x<x]}$. यावरून $\Th{Q} \Proves/
\lforall[x][\lnot x < x]$ हे दिसते.
\end{ex}

\begin{prob}
\ref{mod:mar:mdq:ex:model-K-of-Q} मधील $\Struct{K}$ ही~$\Th{Q}$ ची
उरलेली स्वयंसिद्धके पूर्ण करते हे सिद्ध करा:
\begin{align*}
  & \lforall[x][\eq[(x \times \Obj 0)][\Obj 0]] \tag{$\mathsf{Q}_6$}\\
  & \lforall[x][\lforall[y][\eq[(x \times y')][((x \times y) + x)]]] \tag{$\mathsf{Q}_7$}\\
  & \lforall[x][\lforall[y][(x < y \liff \lexists[z][\eq[(z' + x)][y]])]] \tag{$\mathsf{Q}_8$}
\end{align*}
फक्त~$\prime$ असलेले, $\Struct{N}$ मध्ये सत्य पण $\Struct{K}$ मध्ये असत्य
असे वाक्य शोधा.
\end{prob}

\begin{ex}
\label{mod:mar:mdq:ex:model-L-of-Q} $\Domain{L} = \Nat \cup \{a, b\}$ असा प्रांत
आणि पुढीलप्रमाणे दिलेली $\Assign{\prime}{L} = \nssucc$ व
$\Assign{+}{L} = \nsplus$ ही अर्थनिर्धारणे असलेली संरचना~$\Struct{L}$
विचारात घ्या:
\[\begin{array}{c|c}
  x & x^\nssucc \\
  \hline
  n & n+1 \\
  a & a\\
  b & b
\end{array}
\qquad
\begin{array}{c|ccc}
  x \nsplus y & m & a & b\\
  \hline
  n & n+m & b & a\\
  a & a & b & a\\
  b & b & b & a
\end{array}
\]
$\nssucc$ हे एकास-एक आहे, $0$ त्याच्या व्याप्तीत नाही, आणि
$0$ व्यतिरिक्त प्रत्येक~$x \in \Domain{L}$ त्याच्या व्याप्तीत आहे;
म्हणून $\mathsf{Q}_1$--$\mathsf{Q}_3$ ही स्वयंसिद्धके~$\Struct{L}$ मध्ये सत्य आहेत.
प्रत्येक $x$ साठी $x \nsplus 0 = x$, म्हणून $\mathsf{Q}_4$ सुद्धा सत्य आहे.
$\mathsf{Q}_5$ साठी $x \nsplus y^\nssucc$ आणि $(x \nsplus y)^\nssucc$ विचारात
घ्या. $x$ आणि $y$ दोन्ही मानक असतील, तर ते समान आहेत; कारण तेव्हा
$\nssucc$ आणि $\nsplus$ ही $\prime$ आणि $+$ यांच्याशी जुळतात. $x$ अमानक
आणि $y$ मानक असेल, तर $x \nsplus y^\nssucc = x = x^\nssucc =
(x \nsplus y)^\nssucc$. $x$ आणि $y$ दोन्ही अमानक असतील, तर चार प्रकरणे
आहेत:
\begin{align*}
&  a \nsplus a^\nssucc  = b = b^\nssucc = (a \nsplus a)^\nssucc\\
&  b \nsplus b^\nssucc  = a = a^\nssucc = (b \nsplus b)^\nssucc\\
&  b \nsplus a^\nssucc  = b = b^\nssucc = (b \nsplus a)^\nssucc\\
&  a \nsplus b^\nssucc  = a = a^\nssucc = (a \nsplus b)^\nssucc\\
\intertext{जर $x$ मानक पण $y$ अमानक असेल, तर}
&  n \nsplus a^\nssucc  = n \nsplus a = b = b^\nssucc = (n \nsplus a)^\nssucc\\
&  n \nsplus b^\nssucc  = n \nsplus b = a = a^\nssucc = (n \nsplus b)^\nssucc
\end{align*}

त्यामुळे $\Sat{L}{\mathsf{Q}_5}$. पण $a \nsplus 0 \neq 0 \nsplus a$, म्हणून
$\Sat/{L}{\lforall[x][\lforall[y][\eq[(x+y)][(y+x)]]]}$.
\end{ex}

\begin{prob}
\ref{mod:mar:mdq:ex:model-L-of-Q} मधील $\Struct{L}$ मध्ये
$\times$ आणि $<$ यांचे अर्थनिर्धारण करणारी $\nstimes$ आणि $\nsless$
समाविष्ट करून तिचा विस्तार करा. तुमची रचना~$\Th{Q}$ ची उरलेली
स्वयंसिद्धके पूर्ण करते हे दाखवा:
\begin{align*}
& \lforall[x][\eq[(x \times \Obj 0)][\Obj 0]] \tag{$\mathsf{Q}_6$}\\ &
  \lforall[x][\lforall[y][\eq[(x \times y')][((x \times y) + x)]]]
  \tag{$\mathsf{Q}_7$}\\ & \lforall[x][\lforall[y][(x < y \liff \lexists[z][\eq[(z'+x)][y]])]] \tag{$\mathsf{Q}_8$}
\end{align*}
\end{prob}

\begin{prob}
\ref{mod:mar:mdq:ex:model-L-of-Q} मधील $\Struct{L}$ मध्ये $a^\nssucc
= a$ आणि $b^\nssucc = b$. ज्या~$\Th{Q}$ च्या प्रतिमानात
$a^\nssucc = b$ आणि $b^\nssucc = a$ असे आहे, असे प्रतिमान अस्तित्वात आहे का?
\end{prob}

\begin{explain}
अमानक घटक हे मानक घटकांच्या ``पलीकडे'' असलेली~$\Th{Q}$ ची
प्रतिमाने आपण स्पष्टपणे रचली आहेत. किंबहुना, स्वयंसिद्धकांनुसार तितके असणे
आवश्यक आहे. अमानक घटक~$x$ हा ${} \nsless 0$ असू शकत नाही, कारण
$\Th{Q} \Proves \lforall[x][\lnot x<0]$
(\ref{inc:req:min:lem:less-zero} पाहा). तसेच, प्रत्येक $n$ साठी
$\Th{Q} \Proves \lforall[x][(x < \num{n}' \lif
(\eq[x][\num{0}] \lor \eq[x][\num{1}] \lor \dots \lor
\eq[x][\num{n}]))]$ (\ref{inc:req:min:lem:less-nsucc}); म्हणून
कोणत्याही~$n>0$ साठी $a \nsless n$ असू शकत नाही.
\end{explain}











\section{$\Th{PA}$ ची प्रतिमाने}\label{mod:mar:mpa:sec}

\begin{explain}
$\Th{TA}$ चे कोणतेही अमानक प्रतिमान हे~$\Th{PA}$ चेही अमानक प्रतिमान
असते. $\Th{TA}$ ची आणि म्हणून~$\Th{PA}$ ची अमानक प्रतिमाने अस्तित्वात
आहेत हे आपल्याला माहीत आहे. अशा अमानक प्रतिमानांत अमानक ``संख्या'', म्हणजे
सर्व मानक ``संख्यांच्या'' ``पलीकडे'' असलेले प्रांताचे घटक, असतात
हेही आपल्याला माहीत आहे. पण त्यांची मांडणी कशी असते? त्या किती असतात?
दुर्बलतर उपपत्ती~$\Th{Q}$ च्या प्रतिमानात फक्त एक अमानक संख्या असू शकते
हे आपण पाहिले. पण या साध्या संरचना~$\Th{PA}$ किंवा $\Th{TA}$ ची
प्रतिमाने नाहीत.

$\Th{PA}$ किंवा $\Th{TA}$ च्या प्रतिमानांची रचना समजून घेण्याची गुरुकिल्ली
म्हणजे या उपपत्तींमध्ये कोणती तथ्ये निष्पन्न करता येण्याजोगे आहेत हे पाहणे. उदाहरणार्थ,
$\Th{PA}$ कडून आधीच $\lforall[x][\eq/[x][x']]$ आणि
$\lforall[x][\lforall[y][\eq[(x+y)][(y+x)]]]$ सिद्ध होतात; म्हणून ही
वाक्ये असत्य असलेल्या साध्या रचना~$\Th{PA}$ ची प्रतिमाने असू शकत
नाहीत.

$\Struct{M}$ हे~$\Th{PA}$ चे प्रतिमान आहे असे समजा. मग
$\Th{PA} \Proves \varphi$ असल्यास $\Sat{M}{\varphi}$. पुन्हा
$\Assign{\Obj{0}}{M}$ करिता $\nszero$, $\Assign{\prime}{M}$ करिता
$\nssucc$, $\Assign{+}{M}$ करिता $\nsplus$,
$\Assign{\times}{M}$ करिता $\nstimes$ आणि $\Assign{<}{M}$ करिता
$\nsless$ वापरू. मग कोणतेही वाक्य~$\varphi$ हे $\nszero$, $\nssucc$,
$\nsplus$, $\nstimes$ आणि $\nsless$ यांविषयी एखादी अट सांगते; आणि
$\Sat{M}{\varphi}$ असेल, तर ती अट पूर्ण झाली पाहिजे. उदाहरणार्थ,
$\Sat{M}{\mathsf{Q}_1}$, म्हणजे
$\Sat{M}{\lforall[x][\lforall[y][(\eq[x'][y'] \lif \eq[x][y])]]}$,
असेल तर $\nssucc$ हे एकास-एक असले पाहिजे.
\end{explain}

\begin{prop}
$\Struct{M}$ मध्ये $\nsless$ हा काटेकोर रेषीय क्रम आहे, म्हणजे तो पुढील
अटी पूर्ण करतो:
\begin{enumerate}
\item कोणत्याही~$x \in \Domain{M}$ साठी $x \nsless x$ नाही.
\item $x \nsless y$ आणि $y \nsless z$ असतील, तर $x \nsless z$.
\item कोणत्याही $x \neq y$ साठी $x \nsless y$ किंवा $y \nsless x$.
\end{enumerate}
\end{prop}

\begin{proof}
$\Th{PA}$ कडून पुढील सूत्रे सिद्ध होतात:
\begin{enumerate}
\item $\lforall[x][\lnot x < x]$
\item $\lforall[x][\lforall[y][\lforall[z][((x < y \land y < z) \lif x < z)]]]$
\item $\lforall[x][\lforall[y][((x < y \lor y < x) \lor \eq[x][y]))]]$
\end{enumerate}
\end{proof}

\begin{prop}
\label{mod:mar:mpa:prop:M-discrete} $\nsless$-क्रमात $\nszero$ हा~$\Domain{M}$ मधील
लघुतम घटक आहे. प्रत्येक $x$ साठी $x \nsless x^\nssucc$, आणि हा
गुणधर्म असलेला $\nsless$-लघुतम घटक म्हणजे $x^\nssucc$. शून्याहून
वेगळ्या कोणत्याही~$x$ साठी $y^\nssucc = x$ असा एकमेव $y$ असतो. ($y$ ला
$\Struct{M}$ मधील~$x$ चा ``पूर्ववर्ती'' म्हणतो आणि तो~$^\nssucc x$ असा
दर्शवतो.)

\end{prop}

\begin{proof}
सराव.
\end{proof}

\begin{prob}
\ref{mod:mar:mpa:prop:M-discrete} मधील $\nszero$, $\nssucc$ आणि
$\nsless$ यांचे गुणधर्म सुनिश्चित करणारी, $\Th{PA}$ मध्ये निष्पन्न करता येण्याजोगे
(आणि म्हणून~$\Struct{N}$ मध्ये सत्य) असलेली~$\Lang{L_A}$ मधील
वाक्ये शोधा.
\end{prob}

\begin{prop}
$\Struct{M}$ चे सर्व मानक घटक हे सर्व अमानक घटक पेक्षा
($\nsless$ नुसार) लहान असतात.
\end{prop}

\begin{proof}
$\Struct{M}$ चा मानक घटक~$\Value{\num{n}}{M}$ याचा संक्षेप म्हणून
$n$ वापरू. कोणत्याही~$n \in \Nat$ साठी
$\lforall[x][(x < \num{n}' \lif (\eq[x][\num{0}] \lor
  \eq[x][\num{1}] \lor \dots \lor \eq[x][\num{n}]))]$ हे $\Th{Q}$
कडून आधीच सिद्ध होते. $\nsless \nszero$ असलेले घटक
नाहीत. म्हणून $n$ मानक आणि $x$ अमानक असेल, तर $x \nsless n$ असू शकत
नाही. व्याख्येनुसार, अमानक घटक म्हणजे कोणत्याही~$n \in \Nat$ साठी
$\Value{\num{n}}{M}$ नसलेला घटक; म्हणून $x \neq n$ सुद्धा. $\nsless$
हा रेषीय क्रम असल्यामुळे $n \nsless x$ असलेच पाहिजे.
\end{proof}

\begin{prop}
$\Domain{M}$ मधील प्रत्येक अमानक घटक~$x$ हा पुढील उपसंचाचा घटक
असतो:
\[\dots ^{\nssucc\nssucc\nssucc}x \nsless ^{\nssucc\nssucc}x \nsless
^{\nssucc}x \nsless x \nsless x^{\nssucc} \nsless x^{\nssucc\nssucc}
\nsless x^{\nssucc\nssucc\nssucc} \nsless \dots
\]
या उपसंचाला \emph{$x$ चा खंड} म्हणतो आणि तो $[x]$ असा लिहितो. त्याला
लघुतम किंवा महत्तम घटक नाही. एखाद्या मानक~$n$ साठी
$x \nsplus n = y$ किंवा $y \nsplus n = x$ असे असणाऱ्या
$y \in \Domain{M}$ चा संच, असे त्याचे लक्षण सांगता येते.
\end{prop}

\begin{proof}
असा संच~$[x]$ नेहमी अस्तित्वात असतो हे स्पष्ट आहे; कारण
$\Domain{M}$ मधील प्रत्येक अशून्य घटक~$y$ ला एकमेव
उत्तरवर्ती~$y^\nssucc$ आणि एकमेव पूर्ववर्ती~$^\nssucc y$ असतो. लागोपाठच्या
घटक $y$, $y^\nssucc$ साठी $y \nsless y^\nssucc$, आणि
$y$ ज्याच्यापेक्षा $\nsless$ आहे असा $\Domain{M}$ मधील $\nsless$-लघुतम
घटक म्हणजे $y^\nssucc$. नेहमी $^\nssucc y \nsless y$ आणि
$y \nsless y^\nssucc$ असल्यामुळे $[x]$ ला लघुतम किंवा महत्तम घटक
नाही. $y \in [x]$ असेल, तर $x \in [y]$; कारण मग एकतर
$y^{\nssucc\dots\nssucc} = x$ किंवा $x^{\nssucc\dots\nssucc} = y$.
$y^{\nssucc\dots\nssucc} = x$ ($n$ वेळा $\nssucc$) असेल, तर
$y \nsplus n = x$, आणि याचा व्यत्यासही खरा; कारण $n$ ही $\prime$ ची
संख्या असेल, तर $\Th{PA} \Proves
\lforall[x][\eq[x^{\prime\dots\prime}][(x + \num{n})]]$.
\end{proof}

\begin{prop}
$[x] \neq [y]$ आणि $x \nsless y$ असतील, तर कोणत्याही $u \in [x]$ आणि
कोणत्याही $v \in [y]$ साठी $u \nsless v$.
\end{prop}

\begin{proof}
$\Th{PA} \Proves \lforall[x][\lforall[y][(x < y \lif (x' < y
    \lor x' = y))]]$ हे लक्षात घ्या. म्हणून $u \nsless v$ आणि
$[u] \neq [v]$ असतील, तर कोणत्याही~$n$ साठी
$u \nsplus n^\nssucc \nsless v$ सुद्धा.

प्रत्येक $u \in [x]$ हा~$\nsless y$ आहे: गृहीतकानुसार
$x \nsless y$. $u \nsless x$ असेल, तर संक्रमणशीलतेने $u \nsless y$.
आणि $x \nsless u$ पण $u \in [x]$ असेल, तर कोणत्यातरी~$n$ साठी
$u = x\nsplus n^\nssucc$; म्हणून नुकत्याच सिद्ध केलेल्या तथ्यावरून
$u \nsless y$.

आता $v \in [y]$ हा~$\nsless y$ आहे, म्हणजे एखाद्या मानक~$m$
साठी $v \nsplus m^\nssucc = y$, असे समजा. यामुळे $v \nsless x$ असण्याची
शक्यता बाद होते; अन्यथा $y = v \nsplus m^\nssucc \nsless x$ झाले असते.
तसेच स्पष्टपणे $x \neq v$; अन्यथा
$x \nsplus m^\nssucc = v \nsplus m^\nssucc = y$ आणि $[x] = [y]$ झाले
असते. म्हणून $x \nsless v$. पण मग कोणत्याही~$n$ साठी
$x \nsplus n^\nssucc \nsless v$ सुद्धा. त्यामुळे $x \nsless u$ आणि
$u \in [x]$ असतील, तर $u \nsless v$. $u \nsless x$ असेल, तर
संक्रमणशीलतेने $u \nsless v$.

शेवटी, $y \nsless v$ असेल, तर आपण दाखवल्याप्रमाणे $u \nsless y$ आणि
$y \nsless v$ असल्यामुळे $u \nsless v$.
\end{proof}

\begin{cor}
$[x] \neq [y]$ असेल, तर $[x] \cap [y] = \emptyset$.
\end{cor}

\begin{proof}
$z \in [x]$ आणि $x \nsless y$ असे समजा. मग प्रत्येक $u \in [y]$ साठी
$z \nsless u$. जर $z \in [y]$ असेल, तर $z \nsless z$ मिळेल.
$y \nsless x$ असेल तेव्हाही याचप्रमाणे.
\end{proof}

\begin{explain}
याचा अर्थ खंडांनाच $\nsless$ चा आदर राखणारा क्रम देता येतो:
$[x] \nsless [y]$ हे तेव्हाच, जेव्हा $x \nsless y$; किंवा सममूल्य रीतीने,
कोणत्याही $u \in [x]$ आणि $v \in [y]$ साठी $u \nsless v$. मानक खंड
$[0]$ हा लघुतम खंड आहे हे स्पष्ट आहे. त्याचा कोणत्याही अमानक खंडाशी छेद
नाही, आणि कोणत्याही दोन अमानक खंडांचाही परस्परांशी छेद नाही. विशेषतः,
उत्तरवर्ती किंवा पूर्ववर्ती वारंवार घेऊन वेगळ्या खंडात ``पोहोचता'' येत नाही.
\end{explain}

\begin{prop}
$x$ आणि $y$ अमानक असतील, तर $x \nsless x \nsplus y$ आणि
$x \nsplus y \notin [x]$.
\end{prop}

\begin{proof}
$y$ अमानक असेल, तर $y \neq \nszero$. $\Th{PA} \Proves
\lforall[x][(y \neq \Obj{0} \lif x < (x+y))]$. आता
$x \nsplus y \in [x]$ असे समजा. $x \nsless x \nsplus y$ असल्यामुळे
$x \nsplus n^\nssucc = x \nsplus y$ असे झाले असते. पण $\Th{PA} \Proves
\lforall[x][\lforall[y][\lforall[z][(\eq[(x+y)][(x+z)] \lif y = z)]]]$
(बेरीजेसाठीचा निरसन नियम). याचा अर्थ एखाद्या मानक~$n$ साठी
$y = n^\nssucc$ असा झाला असता; पण $y$ अमानक आहे असे गृहीत आहे.
\end{proof}

\begin{prop}
लघुतम अमानक खंड नाही.
\end{prop}

\begin{proof}
$\Th{PA} \Proves \lforall[x][\lexists[y][(\eq[(y+y)][x] \lor
      \eq[(y+y)'][x])]]$; म्हणजे प्रत्येक $x$ ला~$2$ ने भाग जातो
(शक्यतो बाकी~$1$ ठेवून). $x$ अमानक असेल, तर $y$ सुद्धा अमानक असतो.
आधीच्या विधानावरून $y \nsless y \nsplus y$ आणि
$y \nsplus y \notin [y]$. मग $y \nsless (y \nsplus y)^\nssucc$ आणि
$(y \nsplus y)^\nssucc \notin [y]$ सुद्धा. पण $x = y \nsplus y$ किंवा
$x = (y \nsplus y)^\nssucc$, म्हणून $y \nsless x$ आणि $y \notin [x]$.
\end{proof}

\begin{prop}
महत्तम खंड नाही.
\end{prop}

\begin{proof}
सराव.
\end{proof}

\begin{prob}
$\Th{PA}$ च्या अमानक प्रतिमानात महत्तम खंड नसतो हे दाखवा.
\end{prob}

\begin{prop}
\label{mod:mar:mpa:prop:blocks-dense}
खंडांचा क्रम सघन आहे. म्हणजे, $x \nsless y$ आणि $[x] \neq [y]$ असतील,
तर त्या दोहोंपेक्षा वेगळा आणि त्यांच्यामध्ये असलेला एक खंड~$[z]$ असतो.
\end{prop}

\begin{proof}
$x \nsless y$ असे समजा. पूर्वीप्रमाणे, $x \nsplus y$ ला दोनने भाग जातो
(शक्यतो बाकी ठेवून): एखादा $z \in \Domain{M}$ असा असतो की एकतर
$x \nsplus y = z \nsplus z$ किंवा
$x \nsplus y = (z \nsplus z)^\nssucc$. घटक~$z$ हा $x$ आणि~$y$ यांची
``सरासरी'' आहे, आणि $x \nsless z$ व $z \nsless y$.

\end{proof}

\begin{prob}
\ref{mod:mar:mpa:prop:blocks-dense} चा तपशीलवार पुरावा लिहा.
$z$ चे अस्तित्व सुनिश्चित करण्यासाठी $\Th{PA}$ ने कोणते वाक्य
निष्पन्न केले पाहिजे? $x \nsless z$ आणि $z \nsless y$ का, तसेच
$[x] \neq [z]$ आणि $[z] \neq [y]$ का?
\end{prob}

\begin{explain}
प्रतिमान गणनीय असेल, तर अमानक खंडांचा क्रम परिमेय संख्यांच्या क्रमासारखा
असतो: ते अंतबिंदुविरहित गणनीय अनंत सघन रेषीय क्रम बनवतात. असे
कोणतेही दोन गणनीय अनंत क्रम समरूपी असतात हे दाखवता येते. त्यावरून
सत्य अंकगणिताच्या कोणत्याही दोन गणनीय अमानक प्रतिमान
$\Struct{M}_1$ आणि $\Struct{M_2}$ यांची फक्त $<$ आणि $=$ असलेल्या भाषेतील
न्यूनीकरणे समरूपी आहेत. खरेच, पुढीलप्रमाणे समरूपण~$h$ ची व्याख्या करता
येते: $\Struct{M_1}$ आणि $\Struct{M_2}$ यांचे मानक भाग मानक
प्रतिमान~$\Struct{N}$ शी आणि म्हणून परस्परांशी समरूपी आहेत. अमानक भाग
बनवणाऱ्या खंडांचा क्रम स्वतः परिमेय संख्यांसारखा असल्यामुळे ते समरूपी आहेत;
प्रत्येक खंडातील स्वेच्छ घटक परस्पर जुळवून आणि नंतर $\Struct{M_1}$ मधील
$x$ च्या उत्तरवर्तीची प्रतिमा ही $\Struct{M_2}$ मधील $x$ च्या प्रतिमेचा
उत्तरवर्ती घेऊन, खंडांचे समरूपण खंडांच्या \emph{आत} विस्तारित करता येते.
$\mathfrak{M}_1$ आणि $\mathfrak{M}_2$ या अंकगणिताच्या पूर्ण भाषेत समरूपी
आहेत असे यावरून \emph{निष्पन्न होत नाही} हे लक्षात घ्या (खरे तर समरूपण
नेहमी भाषा सापेक्ष असते), कारण $\Domain{M_1}$ आणि
$\Domain{M_2}$ वर बेरीज आणि गुणाकार परिभाषित करण्याचे परस्पर-असमरूपी मार्ग
आहेत. (पूर्ण उपपत्तीच्या गणनीय प्रतिमानांची संख्या~2 असू शकत नाही, या
वॉट यांच्या प्रसिद्ध प्रमेयावरूनही हे निष्पन्न होते.)

\end{explain}










\section{अंकगणिताची संगणनक्षम प्रतिमाने}\label{mod:mar:cmp:sec}

\begin{explain}
मानक प्रतिमान~$\Struct{N}$ ला दोन उपयुक्त वैशिष्ट्ये आहेत. त्याचा प्रांत
नैसर्गिक संख्यांचा संच~$\Nat$ आहे, म्हणजे त्याचे घटक हे नेमके तेच प्रकारचे
विषय आहेत ज्यांच्याविषयी अंकगणिताची भाषा वापरून आपल्याला बोलायचे असते, आणि
मानक संख्यांक~$\num{n}$ हा खरोखर~$n$ लाच निर्देशित करतो. दुसरे उपयुक्त
वैशिष्ट्य म्हणजे~$\Lang{L_A}$ मधील अतार्किक चिन्हांची सर्व अर्थनिर्धारणे
\emph{संगणनक्षम} आहेत. $\Assign{\prime}{N}$, $\Assign{+}{N}$ आणि
$\Assign{\times}{N}$ म्हणून काम करणारी उत्तरवर्ती, बेरीज आणि गुणाकार ही
फलने संख्यांवरील संगणनक्षम फलने आहेत. (म्हणजे ट्यूरिंग यंत्रांनी संगणन करता
येणारी किंवा, उदाहरणार्थ, आदिम पुनरावर्तनाने परिभाषित करता येणारी.) तसेच
$\Struct{N}$ वरील लघुतर संबंध, म्हणजे~$\Assign{<}{N}$, निर्णेय आहे.

$\Th{Q}$ आणि $\Th{PA}$ यांसारख्या अंकगणितीय उपपत्तींच्या अमानक प्रतिमानांत
अमानक घटक असलेच पाहिजेत. त्यामुळे त्यांच्या प्रांतांत सामान्यतः~$\Nat$
व्यतिरिक्तही घटक असतात. तरी कोणतीही गणनीय संरचना कोणत्याही
गणनीय अनंत संचावर, त्यात~$\Nat$ वरही, उभारता येते. म्हणून प्रांत
$\Nat$ असलेली अमानक प्रतिमानेही आहेत. अशा~$\Struct{M}$ प्रतिमानांत अर्थातच
किमान काही संख्या आपली नेहमीची भूमिका बजावू शकत नाहीत, कारण असा काही~$k$
असतो की प्रत्येक~$n \in \Nat$ साठी तो~$\Value{\num{n}}{M}$ पेक्षा
वेगळा असतो.
\end{explain}

\begin{defn}
$\Lang{L_A}$ साठीची संरचना~$\Struct{M}$ ही \emph{संगणनक्षम}
तेव्हाच, जेव्हा $\Domain{M} = \Nat$ असते, $\Assign{\prime}{M}$,
$\Assign{+}{M}$ आणि $\Assign{\times}{M}$ ही संगणनक्षम फलने असतात, आणि
$\Assign{<}{M}$ हा निर्णेय संबंध असतो.
\end{defn}

\begin{ex}\label{mod:mar:cmp:ex:comp-model-q}
\ref{mod:mar:mdq:ex:model-K-of-Q} मधील रचना~$\Struct{K}$ आठवा. तिचा प्रांत
$\Domain{K} = \Nat
\cup \{a\}$ होता आणि अर्थनिर्धारणे पुढीलप्रमाणे होती:
\begin{align*}
  \Assign{\Obj{0}}{K} & = 0\\
  \Assign{\prime}{K}(x) & =
  \begin{cases}
    x+1 & \text{जर $x\in \Nat$}\\
    a & \text{जर $x = a$}
  \end{cases}\\
  \Assign{+}{K}(x, y) & =
  \begin{cases}
    x+y & \text{जर $x$, $y \in\Nat$}\\
    a & \text{इतर प्रसंगी}
  \end{cases}\\
  \Assign{\times}{K}(x, y) & =
  \begin{cases}
    xy & \text{जर $x$, $y \in\Nat$}\\
    0 & \text{जर $x=0$ किंवा $y=0$}\\
    a & \text{इतर प्रसंगी}\\
  \end{cases}\\
  \Assign{<}{K} & =
  \Setabs{\tuple{x,y}}{x, y \in \Nat \text{ आणि } x<y} \cup
  \Setabs{\tuple{x,a}}{x \in \Domain{K}}
\end{align*}

पण $\Domain{K}$ हा गणनीय अनंत आहे आणि म्हणून तो~$\Nat$ शी तुल्यबल
आहे. उदाहरणार्थ, $g\colon \Nat \to \Domain{K}$ असे घेऊ, जिथे $g(0) =
a$ आणि $n>0$ साठी $g(n) = n-1$; हे एक एकास-एक आच्छादन आहे.

त्याचा वापर करून नव्या~$\Th{Q}$-प्रतिमान~$\Struct{K'}$ आणि
$\Struct{K}$ यांमध्ये समरूपता मिळवता येते. $\Struct{K'}$ मध्ये
$\Lang{L_A}$ च्या चिन्हांना वेगळी फलने आणि संबंध नेमावे लागतात, कारण
$\Nat$ मधील वेगवेगळे घटक मानक आणि अमानक संख्यांच्या भूमिका
बजावतात.

विशेषतः, आता~$0$ हा लघुतम मानक संख्या नसून~$a$ ची भूमिका बजावतो. आता~$1$
ही लघुतम मानक संख्या आहे. म्हणून $\Assign{\Obj{0}}{K'} = 1$ असे नेमू.
आता उत्तरवर्ती फलनही वेगळे आहे: एखादी मानक संख्या, म्हणजे $n > 0$, दिली
असता ते अजूनही $n+1$ परत करते. पण आता~$0$ हा~$a$ ची भूमिका बजावतो आणि
तो स्वतःचाच उत्तरवर्ती आहे. म्हणून $\Assign{\prime}{K'}(0) = 0$.
बेरीज आणि गुणाकारासाठी याचप्रमाणे
\begin{align*}
\Assign{+}{K'}(x, y) & =
  \begin{cases}
    x+y-1 & \text{जर $x$, $y >0$}\\
    0 & \text{इतर प्रसंगी}
  \end{cases}\\
  \Assign{\times}{K'}(x, y) & =
  \begin{cases}
    1 & \text{जर $x = 1$ किंवा $y = 1$}\\
    xy - x - y + 2 & \text{जर $x$, $y > 1$}\\
    0 & \text{इतर प्रसंगी}\\
  \end{cases}
\end{align*}
आणि $\tuple{x, y} \in \Assign{<}{K'}$ हे तेव्हाच, जेव्हा $x < y$, $x > 0$
आणि $y > 0$ असतात, किंवा $y = 0$ असतो.

ही सर्व फलने नैसर्गिक संख्यांची संगणनक्षम फलने आहेत आणि
$\Assign{<}{K'}$ हा~$\Nat$ वरील निर्णेय संबंध आहे---पण ती~$\Nat$ वरील
उत्तरवर्ती, बेरीज आणि गुणाकार हीच फलने नाहीत, तसेच $\Assign{<}{K'}$ हाही
$\Nat$ वरील~$<$ हाच संबंध नाही.
\end{ex}

\begin{prob}
\ref{mod:mar:mdq:ex:model-L-of-Q} मधील~$\Struct{L}$ शी समरूपी,
$\Domain{L'} = \Nat$ असलेली संरचना~$\Struct{L'}$ द्या.
\end{prob}

\begin{explain}
\ref{mod:mar:cmp:ex:comp-model-q} यावरून~$\Th{Q}$ ला प्रांत~$\Nat$ असलेली
संगणनक्षम अमानक प्रतिमाने आहेत हे दिसते. पण पुढील निष्कर्ष दाखवतो की
$\Th{PA}$ च्या प्रतिमानांसाठी (आणि म्हणून~$\Th{TA}$ च्या प्रतिमानांसाठीही)
हे खरे नाही.
\end{explain}

\begin{thm}[टेनेनबाउमचे प्रमेय]
$\Struct{N}$ हे~$\Th{PA}$ चे एकमेव संगणनक्षम प्रतिमान आहे.
\end{thm}



